\documentclass[10pt,a4paper]{amsart}
\usepackage[margin=19mm]{geometry}
\usepackage{amsmath,amssymb,mathtools}
\usepackage[scr=rsfs,cal=euler]{mathalfa}
\usepackage{xfrac}
\usepackage[unicode,hidelinks,bookmarks=false]{hyperref}
\newcommand{\ie}{i.e.\ }
\newcommand{\cf}{cf.\ }
\newcommand{\resp}{respectively\ }
\newcommand{\Subsection}{\subsection}
\providecommand{\nobreakdash}{\nobreak\hbox{-}}
\setlength{\emergencystretch}{2em}
\multlinegap0pt
\usepackage{bm}
\usepackage{bbm}
\usepackage{dsfont}
\usepackage{xspace}
\usepackage{cancel}
\usepackage{mathtools}
\usepackage{amsthm}
\makeatletter
\@ifundefined{theorem}{\newtheorem{theorem}{Theorem}}{}
\@ifundefined{lemma}{\newtheorem{lemma}[theorem]{Lemma}}{}
\makeatother
\newcommand{\RR}{\mathbb{R}}
\title[Distinct permutation dot products]{Distinct permutation dot products}
\author{Cosmin Pohoata}
\date{Source: arXiv:2601.12445v1 (CC BY 4.0)}
\begin{document}
\maketitle

\noindent\emph{Source and licence.} Distinct permutation dot products. Version arXiv:2601.12445v1 (CC BY 4.0), distributed under the Creative Commons Attribution 4.0 International licence, as recorded in the arXiv licence metadata for this exact version and shown on the abstract page https://arxiv.org/abs/2601.12445v1. Full rights notice: licenses/arxiv-2601.12445-CC-BY-4.0.txt.\par\smallskip

\noindent\textit{Editorial scope.} Counted unit: the Lemma whose source label is \texttt{lem:two-area-warmup} in the article (source0.tex lines 485--490, source label \texttt{lem:two-area-warmup}), delivered together with the article's own complete original proof of it (source0.tex lines 547--584, one \texttt{proof} environment of 38 lines). No mathematics is written, repaired or re-derived: statement and proof are the article's own text line for line. The proof is detached from the statement in the source: the article states the result at lines 485--490 and prints its proof at lines 547--584, and the proof environment's own header names the statement, so the association is the article's own. Required context retained uncounted: none. Every definition, display and label of those lines is retained unchanged. Omitted from this package and not redistributed: 1--484, 491--546, 585--803. Nothing inside a retained excerpt is omitted or altered: every retained body is delivered exactly as the article's lines are written, so no omission receipt of any kind accompanies this package. Citations: the retained excerpts carry no citation command at all, so no bibliography entry of the article is transcribed and no reference is redistributed. Reference adaptation: none was needed and none was made; every reference inside the delivered unit resolves inside it, so no unresolved reference or citation is produced. Printed slips kept literal: no printed word, symbol, hypothesis or number of the counted statement or of its proof was altered. Layout and class: the article's own class is \texttt{amsart} with packages including fontenc, lmodern, microtype, geometry, amsmath, amssymb, amsthm, mathtools, enumitem, hyperref, cleveref; the wrapper is the lane's pinned \texttt{amsart} editorial wrapper at 10pt on A4, which restates the theorem-like environments the retained text needs and adds the article's own macro definitions that the retained text needs (\texttt{\textbackslash RR}), with extra wrapper packages (bm, bbm, dsfont, xspace, cancel, mathtools). The delivered numbering is the unit's own. Assets: no retained span contains a figure environment, a graphics-inclusion command, a PDF-page inclusion, a table or a TikZ picture; no image file is copied into this package, the package declares no asset dependency and no image file is redistributed with it. Licence: version arXiv:2601.12445v1 of this article is distributed under the Creative Commons Attribution 4.0 International licence, as recorded in the arXiv licence metadata for this exact version and shown on the abstract page https://arxiv.org/abs/2601.12445v1. A full rights notice is delivered at licenses/arxiv-2601.12445-CC-BY-4.0.txt. Characters: every retained excerpt is pure ASCII, so the delivered engine, XeLaTeX with its default Latin Modern text font, reports no missing character.
\begin{lemma}\label{lem:two-area-warmup}
Let $A,B\subset\RR$ be finite sets with $|A|,|B| \geq 2$. Then
\[
\bigl|(A-A)(B-B)+(A-A)(B-B)\bigr|\ \ge\ \frac{|A||B|}{2}.
\]
\end{lemma}

\begin{proof}
Like in the proof of Lemma \ref{lem:two-area-warmup}, write $A=\{a_1<\cdots<a_n\}$ and $B=\{b_1<\cdots<b_m\}$, and define 
\[
\delta:=\min_{1\le t\le n-1}(a_{t+1}-a_t)>0,
\qquad
M:=b_m-b_1>0.
\]
Set
\[
P:=\{\,b_i-b_2:\ 3\le i\le m-1\,\}\subseteq B-B,
\]
so $|P|=m-3$, and each $x\in P$ is represented by a difference $b_i-b_2$ using indices in
$\{2,3,\dots,m-1\}$, hence disjoint from $\{1,m\}$.

Choose an index $s \in [n-1]$ with $a_{s+1}-a_s=\delta$, and select an anchor $a_\star\in A\setminus\{a_s,a_{s+1}\}$. For concreteness, take $a_\star=a_1$ if
$a_1\notin\{a_s,a_{s+1}\}$ and otherwise take $a_\star=a_n$.  Let $I$ be the set of odd indices
$2j-1$ with $1\le j\le\lceil n/2\rceil$, and delete from $I$ any indices corresponding to the elements $\{a_s,a_{s+1}\}$ and also delete the anchor index itself.
Define
\[
P':=\{\,a_t-a_\star:\ t\in I\,\}\subseteq A-A.
\]
Then $|P'|\ge \lceil n/2\rceil-2$.  Moreover, any two distinct elements of $P'$ differ by at least
$2\delta$. Like in the proof of Lemma \ref{lem:two-area-warmup}, consider
\[
U:=\delta P+MP'=\{\delta x+My:\ x\in P,\ y\in P'\} \subset (A-A)(B-B)+(A-A)(B-B).
\]
The map $(x,y)\mapsto \delta x+My$ is still an injective map on $P\times P'$, therefore 
$$|U| \geq |P||P'|\ge (m-3)(\lceil n/2\rceil-2).$$

Most importantly, note that every $u=\delta x+My\in U$ admits a representation using disjoint indices: write
$x=b_i-b_2$ with $3\le i\le m-1$ and write $y=a_t-a_\star$ with $a_t\notin\{a_s,a_{s+1}\}$ and
$a_\star\notin\{a_s,a_{s+1}\}$.  Then
\[
u=(a_{s+1}-a_s)(b_i-b_2)\;+\;(a_t-a_\star)(b_m-b_1),
\]
and the $A$--indices $\{s,s+1,t,\star\}$ are distinct by construction, while the $B$--indices
$\{i,2,m,1\}$ are distinct.  Taking $U(A,B):=U$ completes the proof.
\end{proof}

\end{document}
